% ============================================================================
% Fichier  : partie4/P04_C11_supprimes_artefacts.tex : FORN (LPI342)
% Sources : ISO/IEC 27042 ; NIST SP 800-86. Opération TSANAGA.
% ============================================================================

\chapter{Fichiers Supprimés, Fragments et Artefacts}
\label{chap:artefacts}
\chaptermark{Supprimés et artefacts}

\epigraph{%
    « Supprimer, c'est souvent seulement cacher. »%
}{Principe du cours}

\niveauChapitre{Intermédiaire à expert}{Chapitre~\ref{chap:disque}.}

\begin{boxobjectifs}
\begin{itemize}
    \item Comprendre pourquoi un fichier supprimé subsiste.
    \item Récupérer un fichier par ses fragments.
    \item Identifier un type de fichier par sa signature.
    \item Lire les artefacts d'usage (récents, raccourcis, corbeille).
    \item Qualifier un artefact récupéré.
\end{itemize}
\end{boxobjectifs}

\bigskip

% ============================================================================
\section*{Question directrice : le fichier clients, supprimé par M. NDONGO, peut-il encore être retrouvé ?}
% ============================================================================

M.~NDONGO a supprimé le fichier et vidé la corbeille. Pourtant, les données
peuvent subsister. Ce chapitre explique pourquoi, et comment les récupérer et
les qualifier.

\medskip
\textbf{Plan du chapitre.}
\begin{enumerate}
    \item Pourquoi un fichier subsiste (section~\ref{sec:c11-subsiste}).
    \item Récupérer par fragments (section~\ref{sec:c11-carving}).
    \item Les signatures de fichiers (section~\ref{sec:c11-signatures}).
    \item Les artefacts d'usage (section~\ref{sec:c11-usage}).
\end{enumerate}

\separateur

% ============================================================================
\section{Pourquoi un Fichier Subsiste}
\label{sec:c11-subsiste}
% ============================================================================

Supprimer un fichier ne l'efface pas~: le système marque son espace comme
réutilisable et retire son entrée de l'index, mais les données restent tant
qu'elles ne sont pas réécrites. D'où la possibilité de les
\termecle{récupérer}, en partie ou en totalité, jusqu'à leur écrasement.

% ============================================================================
\section{Récupérer par Fragments}
\label{sec:c11-carving}
% ============================================================================

\begin{boxdef}
\termecle{Récupération par signatures} (\textit{file carving})~: technique qui
reconstitue des fichiers à partir de leurs motifs de début et de fin dans
l'espace non alloué, sans s'appuyer sur l'index du système de fichiers. Elle
retrouve des fichiers dont l'entrée a disparu.
\end{boxdef}

% ============================================================================
\section{Les Signatures de Fichiers}
\label{sec:c11-signatures}
% ============================================================================

Un fichier commence souvent par une \termecle{signature} (quelques octets)
qui trahit son type, indépendamment de son extension. Un fichier renommé en
\texttt{.txt} mais commençant par la signature d'une archive reste une
archive~: la signature prime sur l'extension.

\begin{boxalert}
\textbf{Violation typique.} Se fier à l'extension d'un fichier.

\smallskip
\textbf{Règle~:} l'extension se change en un clic. On identifie le type par la
signature, pas par le nom~: c'est une parade classique de dissimulation.
\end{boxalert}

% ============================================================================
\section{Les Artefacts d'Usage}
\label{sec:c11-usage}
% ============================================================================

\begin{tableaucours}{L{3.6cm}L{8.6cm}}{%
    \tableauHeader{Artefact} &
    \tableauHeader{Ce qu'il révèle}
}
Fichiers récents, raccourcis & Documents ouverts, chemins, volumes amovibles \\
Corbeille & Fichiers supprimés récupérables, date \\
Journaux d'événements & Branchement de périphériques, sessions \\
Fichiers temporaires, caches & Copies de travail, miniatures \\
\end{tableaucours}
\vspace{-0.8em}
\begin{center}{\small\itshape Tableau~11.1. Artefacts d'usage utiles}\end{center}

\begin{boxscenario}{Opération TSANAGA. Le fichier retrouvé}
\textbf{Constatation.} Dans l'espace non alloué, la récupération par
signatures reconstitue un tableur dont la signature est celle d'un classeur,
bien que l'entrée ait disparu. Son contenu correspond au fichier clients. Un
raccourci récent pointe vers un volume amovible, le 3~janvier.

\medskip
\textbf{Analyse} (confiance élevée pour l'existence, moyenne pour l'auteur).
Le fichier supprimé est récupéré et son contenu établi~; la copie vers
l'amovible est corroborée par l'artefact de raccourci. Le faisceau se
renforce.
\end{boxscenario}

\begin{boxobjet}{Fichier récupéré : clients (classeur)}
\textbf{Nature~:} classeur récupéré par signatures dans l'espace non alloué.
\textbf{Source~:} image disque TSA-2026-01 (condensat au dossier).
\textbf{Pertinence~:} contenu correspondant au fichier clients.
\textbf{Fiabilité~:} intègre (travail sur l'image). \textbf{Opposabilité~:}
chaîne de possession continue.
\end{boxobjet}

\begin{boxmetier}{Analyste forensique}
\textbf{Ne pas se fier au nom.} L'analyste identifie les fichiers par leur
signature, fouille l'espace non alloué, et recoupe les artefacts d'usage. Ce
que l'utilisateur croyait caché resurgit souvent.
\end{boxmetier}

\separateur

\begin{boxresume}
\textbf{Ce qu'il faut retenir de ce chapitre~:}
\begin{itemize}
    \item \textbf{Supprimer n'est pas effacer}~: récupération possible.
    \item \textbf{La récupération par signatures} reconstitue sans l'index.
    \item \textbf{La signature} prime sur l'extension.
    \item \textbf{Les artefacts d'usage} révèlent documents et périphériques.
    \item \textbf{On qualifie} un fichier récupéré.
\end{itemize}
\end{boxresume}

\begin{boxinfo}
\textbf{Points de contrôle de l'élève.} Je sais~: expliquer la subsistance
d'un fichier supprimé~; décrire la récupération par signatures~; identifier un
type par signature~; lire les artefacts d'usage~; qualifier un fichier
récupéré.
\end{boxinfo}

\bigskip

\begin{boxexo}[title={\ding{45}\quad Exercice~11.1. Subsister \niveaubase}]
Pourquoi un fichier supprimé peut-il être récupéré~?
\end{boxexo}

\begin{boxexo}[title={\ding{45}\quad Exercice~11.2. Carving \niveaubase}]
Qu'est-ce que la récupération par signatures~?
\end{boxexo}

\begin{boxexo}[title={\ding{45}\quad Exercice~11.3. Extension \niveaubase}]
Pourquoi ne pas se fier à l'extension d'un fichier~?
\end{boxexo}

\begin{boxexo}[title={\ding{45}\quad Exercice~11.4. Signature \niveaumoyen}]
Comment identifier un fichier renommé pour dissimuler son type~?
\end{boxexo}

\begin{boxexo}[title={\ding{45}\quad Exercice~11.5. TSANAGA \niveaumoyen}]
Comment le fichier clients supprimé a-t-il été retrouvé et corroboré~?
\end{boxexo}

\begin{boxexo}[title={\ding{45}\quad Exercice~11.6. Raccourci \niveaumoyen}]
Que révèle un raccourci récent pointant vers un volume amovible~?
\end{boxexo}

\begin{boxexo}[title={\ding{45}\quad Exercice~11.7. Écrasement \niveaumoyen}]
Qu'est-ce qui limite la récupération d'un fichier supprimé~?
\end{boxexo}

\begin{boxexo}[title={\ding{45}\quad Exercice~11.8. Corbeille \niveaumoyen}]
Que peut révéler la corbeille, même vidée~?
\end{boxexo}

\begin{boxexo}[title={\ding{45}\quad Exercice~11.9. Faisceau \niveauexpert}]
Construisez le faisceau d'indices de la copie du fichier clients.
\end{boxexo}

\begin{boxexo}[title={\ding{45}\quad Exercice~11.10. Qualification \niveauexpert}]
Qualifiez complètement le fichier récupéré (cinq rubriques).
\end{boxexo}

\bigskip

\begin{boxinfo}
\textbf{Pour aller plus loin.}
\begin{itemize}
    \item ISO/IEC 27042~\cite{iso27042}~; NIST SP~800-86~\cite{nist80086}.
\end{itemize}
\end{boxinfo}

\begin{boxinfo}
\textbf{Transition.} Le disque ne dit pas tout~: la mémoire garde l'état
vivant. Le Chapitre~\ref{chap:memoire} traite du forensique mémoire.
\end{boxinfo}
